% ---- proof of Theorem T11 (Section 7)
\paragraph{Proof of Theorem~\ref{thm:T11}.}
Substituting~\eqref{eq:def_dH} and~\eqref{eq:def_Q} into~\eqref{eq:kl_pairwise} yields
\[
D_{12}(u) = \tfrac{1}{2} \langle u, Q u \rangle.
\]
The constrained supremum \(\sup_{\|u\|^2 \le E} \langle u, Q u\rangle\) is a Rayleigh quotient problem. Because \(Q \succeq 0\), the Rayleigh--Ritz theorem gives
\[
\sup_{\|u\|^2 \le E} \langle u, Q u\rangle = E\,\lambda_{\max}(Q),
\]
with equality at any unit principal eigenvector whenever the top spectral value is an eigenvalue. This is automatic in the finite-dimensional discretization used for computation. In the infinite-dimensional input-function setting, attainment additionally requires the whitened difference operator \(A = R^{-1/2}\Delta H : \mathcal{U} \to \mathcal{Y}\) to be compact between Hilbert spaces: then \(Q = A^{*}A\) is compact, self-adjoint, and positive, its spectral radius is an eigenvalue, and a principal eigenfunction exists. Rescaling by \(\sqrt{E}\) saturates the energy constraint and attains the bound. Without compactness or an explicit finite discretization, the supremum value~\eqref{eq:T11_max} still holds but need not be attained, and no eigenfunction maximizer should be claimed.

% ---- proof of Theorem T12
\paragraph{Proof of Theorem~\ref{thm:T12}.}
\emph{Existence.} The feasible set of nonnegative measures with total mass at most \(E\) is weakly compact, and the constraints define a closed affine slice, so the supremum is attained.

\emph{Finite support.} Take any nonzero feasible \(\xi\) and normalize it to a probability measure \(\tilde\xi = \xi / \xi(\Omega)\). The vector of scaled pairwise scores
\[
\bar w(\tilde\xi) :=
\left( \int_{\Omega} w_1\, d\tilde\xi, \dots, \int_{\Omega} w_P\, d\tilde\xi \right)
\in \mathbb{R}^P
\]
lies in the convex hull of the curve
\[
\mathcal{W} := \big\{ (w_1(\omega), \dots, w_P(\omega)) : \omega \in \Omega \big\} \subset \mathbb{R}^P.
\]
By Carath\'eodory's theorem in \(\mathbb{R}^P\), every point in \(\operatorname{conv}(\mathcal{W})\) admits a convex representation using at most \(P+1\) points of \(\mathcal{W}\). Therefore \(\tilde\xi\) (and hence \(\xi\)) may be replaced by a measure supported on at most \(P+1\) frequencies without changing any of the \(P\) integrated scores. Rescaling back to total mass \(E\) preserves feasibility and the objective value \(t\).

% ---- practical observation design / disc-vs-param / constrained implementation
\subsection{Practical observation design: spacing, cost, and channel selection}
\label{sec:practical_obs}

The top-\(n\) rule of Theorem~\ref{thm:T13} is exact but idealized. Real experiments impose three corrections.

\subsubsection*{Minimum-spacing constraint}

For a single observation channel with candidate times \(t_1 < \cdots < t_N\) and a minimum index spacing \(L\) between consecutive observations, the constrained exact algorithm is the dynamic program
\begin{equation}
V(k, j) = \max\bigl\{ V(k-1, j),\;\; s(t_k) + V(k - L, j - 1) \bigr\},
\label{eq:dp_spacing}
\end{equation}
where \(V(k, j)\) is the best achievable total score using \(j\) observations among the first \(k\) candidate times. This solves the spacing-constrained pairwise schedule in \(O(N n)\) operations.

\subsubsection*{Variable costs and budget}

When each candidate \((t, o)\) carries a measurement cost \(c(t, o)\) and the total sampling budget is \(B\), select binary variables \(z_{t,o} \in \{0, 1\}\) and solve the 0--1 knapsack
\begin{equation}
\max \sum_{t, o} s_{12}(t, o)\, z_{t, o}
\qquad\text{s.t.}\qquad
\sum_{t, o} c(t, o)\, z_{t, o} \le B,\quad z_{t, o} \in \{0, 1\}.
\label{eq:knapsack}
\end{equation}
For \(P\) model pairs, introduce a worst-case floor \(q\) and the \(P\) constraints \(\sum_{t, o} s_p(t, o)\, z_{t, o} \ge q\) for every pair \(p\), yielding a maximin mixed-integer linear program. This is the observation-side analogue of~\eqref{eq:robust_LP}.

\subsubsection*{Channel selection from the strong-Allee linearization}

The observability analysis of Section~\ref{sec:linearization} establishes three structural facts that frame channel choice:
(i) prey-only observation is structurally sufficient;
(ii) predator-only observation is structurally sufficient;
(iii) observing both species is \emph{not} mathematically mandatory for state observability.

However, the direct-channel asymptotics of Theorem~\ref{thm:T4} refine this picture for discrimination:
\begin{itemize}
\item Perturb prey and measure prey to expose the leading \(s^{-\alpha}\) signature directly.
\item Perturb predator and measure predator for the same reason via the predator self-channel.
\item Use the cross-species channel (prey input, predator output, or vice versa) to estimate the interaction parameters \(c, d\). Cross channels have \(CB = 0\), so their leading asymptotic decay is \(s^{-2\alpha}\) (or lower) with a correspondingly different high-frequency phase; the collocated phase gap \((1-\alpha)\pi/2\) of~\eqref{eq:phase} must not be attributed to them.
\end{itemize}
The conclusion is that the preferred practical design is \emph{not} ``both species at every time''. It is:
\begin{enumerate}
\item high-rate collocated measurement around perturbation/switching times, where the principal eigenfunction of Theorem~\ref{thm:T11} concentrates energy;
\item lower-rate cross-species measurement across the recovery interval, to resolve interaction parameters;
\item environmental covariate measurement only when the Popov--Belevitch--Hautus latent-mode condition is weak or prior uncertainty makes latent forcing competitive.
\end{enumerate}

\subsection{Why discrimination and parameter-estimation designs differ}
\label{sec:disc_vs_param}

It remains to state precisely why the schedules of Theorems~\ref{thm:T11}--\ref{thm:T13} are not interchangeable with designs aimed at parameter estimation. For a parameter vector \(\psi\), local parameter estimation uses the sensitivity operator
\begin{equation}
S_\psi = \frac{\partial H(\psi)}{\partial \psi},
\label{eq:sens_op}
\end{equation}
evaluated at a nominal model. The associated Fisher information is
\begin{equation}
F(u) = S_\psi(u)^{*} R^{-1} S_\psi(u),
\label{eq:fisher}
\end{equation}
whereas the discrimination objective uses the inter-model operator
\begin{equation}
Q_{12} = \Delta H^{*} R^{-1} \Delta H.
\label{eq:Q12_again}
\end{equation}
These are different operators. A design that makes \(F\) well-conditioned can simultaneously leave \(\Delta H u\) small, because a parameter-tunable latent model tracks the fractional response at the nominal operating point --- exactly the failure mode classified as an adaptive-model risk in Section~\ref{sec:separation}. Conversely, a design that maximizes pairwise divergence maximin-robustly (Theorem~\ref{thm:T12}) need not render every parameter of either model individually identifiable. D-optimal parameter design maximizes \(\log\det F\), an objective that rewards complementary sensitivity directions; the top-\(n\) separation theorem (Theorem~\ref{thm:T13}) does \emph{not} apply to D-optimal schedules because the determinant is not additive.

This operational distinction is one of the central findings of the paper: \emph{good inputs for parameter estimation are not automatically good inputs for model discrimination}, and vice versa. The experiment protocol developed in later sections therefore separates the two objectives and runs a two-stage design: a discrimination stage that maximizes the worst-case pairwise divergence via Theorems~\ref{thm:T11}--\ref{thm:T12}, followed by an estimation stage conditioned on the selected model, planned via~\eqref{eq:fisher} with the discrimination data as prior.

\subsection{Implementation under additional constraints}
\label{sec:constrained_impl}

The unconstrained eigenvector solution of Theorem~\ref{thm:T11} is the benchmark and the upper bound for any constrained implementation. When box, rate, or cumulative-impact constraints apply, the design becomes
\begin{equation}
\max_u\; \tfrac{1}{2}\, u^{*} Q u
\qquad\text{s.t.}\qquad
|u_k| \le u_{\max},\quad
|u_{k+1} - u_k| \le r_{\max},\quad
\sum_k |u_k|\, \Delta t \le B,
\label{eq:constrained_opt}
\end{equation}
which is generally a nonconvex quadratic maximization. Reliable solution methods include:
\begin{itemize}
\item direct transcription with multiple random starts, exploiting the structure of \(Q\);
\item semidefinite relaxation with randomized rounding, certifying the eigenvector optimum as the outer bound;
\item dynamic programming for low-dimensional discrete input alphabets;
\item Bayesian optimization when the forward ecological simulator is expensive and gradients of \(H_i\) are unavailable.
\end{itemize}
In every case the unconstrained \(u^{*}\) of~\eqref{eq:T11_input} remains the natural warm start and the ceiling against which realized performance is scored. The finite-support guarantee of Theorem~\ref{thm:T12} further reduces the search space: once the \(P+1\) optimal tones are identified, the constrained problem reduces to choosing their complex amplitudes subject to the physical constraints, a much smaller nonconvex program than full-waveform input optimization.

The combined input-plus-observation design of this section --- Theorems~\ref{thm:T11}, \ref{thm:T12}, and~\ref{thm:T13} together with the practical refinements of Section~\ref{sec:practical_obs} --- is what converts the structural separation of Section~\ref{sec:separation} and the complexity-bounded approximation barrier of Section~\ref{sec:approximation} into a deployable, computable discrimination experiment. The remaining sections address the safety of that experiment (Section~\ref{sec:safety}), its Bayesian extension beyond the linear-Gaussian regime (Section~\ref{sec:bayesian}), and a reproducible benchmark (Section~\ref{sec:benchmark}) that implements the constructions above on the strong-Allee predator--prey model.
